\documentclass{rapportClass}

\usepackage{booktabs}
\usepackage{tabularx}
\usepackage{enumitem}
\setlist{nosep, leftmargin=*}

% Sections numérotées 1, 2, 3 (pas de chapitres : chaque chapitre ouvrirait une nouvelle page).
\renewcommand{\thesection}{\arabic{section}}
\renewcommand{\sectionmark}[1]{\markboth{#1}{}}
\fancyhead[C]{M2 SLCD}
\fancyhead[R]{WeatherScope}

% Parties à compléter par l'équipe distribuée / mesures.
\newcommand{\afaire}[1]{\textcolor{red}{\textbf{[À compléter : #1]}}}
\newcommand{\cmd}[1]{\texttt{#1}}

\title{WeatherScope}
\author{Hugo CLAVEILLE, JULES Stevenson, DIALLO Mamadou Ougailou}
\date{8 novembre 2026}

\fancypagestyle{plain}{
  \fancyhf{}
  \fancyhead[L]{Rapport de Projet}
  \fancyhead[C]{M2 SLCD}
  \fancyhead[R]{BigData \& noSQL}
  \fancyfoot[C]{\thepage}
  \renewcommand{\headrulewidth}{0.4pt}
}

\begin{document}

\begin{titlepage}
    \thispagestyle{empty}
    \begin{center}
        \includegraphics[width=1\textwidth]{assets/logo.png}\\[2cm]
        {\Huge \textbf{WeatherScope}}\\[0.5cm]
        {\Large Observatoire météo historique sur NOAA GHCN-Daily}\\[1cm]
        {\Large BigData \& noSQL}\\[1cm]
        \vspace{1cm}
        \textbf{Hugo CLAVEILLE\\ JULES Stevenson\\ DIALLO Mamadou Ougailou}\\[0.5cm]
        {\Large Master 2 SLCD}\\[2cm]
        \vfill
        {\large 8 novembre 2026}
    \end{center}
\end{titlepage}

\setcounter{page}{1}

% ======================================================================================================================
\section{Périmètre}

\begin{itemize}
  \item \textbf{Zones :} Allemagne (\cmd{GM}) et Espagne (\cmd{SP}), définies par le préfixe pays de l'identifiant NOAA.
    Deux climats contrastés (océanique/continental et méditerranéen), différents des zones des autres groupes (États-Unis, Afrique).
  \item \textbf{Années d'étude :} 2015--2024 (10 ans). \textbf{Référence :} 1991--2020 (normale climatologique de 30 ans),
    chargée par le même pipeline et comptée à part.
  \item \textbf{Stations :} les 774 stations (589 GM, 185 SP) dont au moins un des éléments TMAX, TMIN, PRCP
    est actif pendant 2015--2024 d'après \cmd{ghcnd-inventory.txt}. Les stations qui commencent ou s'arrêtent
    pendant la période sont gardées et signalées (annexe~\ref{ann:anomalies}).
  \item Les seuils du sujet sont dépassés : 2\,650\,999 couples station-jour sur 2015--2024 (minimum 200\,000)
    et 21,46 millions d'observations au total (minimum 10 millions pour E1).
\end{itemize}

% ======================================================================================================================
\section{Données}

\paragraph{Provenance.} Une archive NOAA par station (\cmd{by\_station/<id>.csv.gz}, 774 fichiers, 237 Mo),
récupérée le 8 novembre 2026, plus les fichiers de métadonnées (\cmd{ghcnd-stations.txt}, \cmd{ghcnd-inventory.txt}).
L'URL, la date, la taille et le SHA-256 de chaque fichier sont dans \cmd{data/manifest/}. En base, chaque import
crée un lot (table \cmd{batches}) et chaque observation garde l'identifiant du lot qui l'a écrite.

\paragraph{Volume.} Calculé directement sur les fichiers NOAA (\cmd{verify/data\_report.py}).
Valide = sans \cmd{qflag} et sans \cmd{mflag} P.

\begin{table}[H]
\centering\small
\begin{tabular}{llrrrr}
\toprule
Zone & Élément & Stations 2015--24 & Valides 2015--24 & Valides 1991--2014 (réf.) & Invalides \\
\midrule
GM & TMAX & 474 & 1\,662\,847 & 3\,572\,600 & 104 \\
GM & TMIN & 474 & 1\,662\,826 & 3\,571\,251 & 137 \\
GM & PRCP & 579 & 2\,010\,884 & 4\,208\,967 & 49 \\
SP & TMAX & 183 & 598\,648 & 1\,002\,404 & 622 \\
SP & TMIN & 183 & 603\,741 & 994\,030 & 477 \\
SP & PRCP & 184 & 604\,505 & 965\,962 & 14 \\
\bottomrule
\end{tabular}
\end{table}

\paragraph{Couverture.} Couples station-jour 2015--2024 : 2\,029\,108 (GM), 621\,891 (SP).
Répartition géographique : 213 mailles de 0,5° occupées en Allemagne, 94 en Espagne ; la densité est
plus forte en Allemagne et inégale à l'intérieur de chaque pays (côtes et grandes villes en Espagne).
Le volume est stable de 1991 à 2023 ; en 2024, l'Espagne perd environ 24\,\% de ses observations
(142\,937 contre 187\,782 en 2023), car des stations s'arrêtent ou ne sont pas encore publiées.

\paragraph{Traitement des données.}
\begin{itemize}
  \item \textbf{\cmd{qflag} non vide :} observation stockée et comptée, mais exclue de tous les calculs (1\,403 lignes).
  \item \textbf{Pluie T (trace) :} valide, comptée 0~mm et à part (267 lignes, toutes en Espagne). Une trace est une
    vraie observation d'une pluie négligeable.
  \item \textbf{Pluie P (manquante, présumée nulle) :} traitée comme manquante, jamais comme un zéro mesuré
    (aucune dans notre extrait, règle tout de même implémentée).
  \item \textbf{Pluie B, D :} cumuls de 2 ou 4 relevés partiels, ce sont de vraies mesures, gardées telles quelles.
  \item \textbf{Valeurs inhabituelles sans \cmd{qflag} :} gardées (ex. 203,7~mm le 13/09/2019 à SP000007038,
    épisode méditerranéen réel) ; seule la NOAA invalide une valeur.
  \item \textbf{Historiques incomplets :} les jours manquants ne sont jamais complétés. Ils coupent les épisodes (F2),
    rendent un mois inéligible sous 80\,\% de jours valides (F3) et sont pris en compte dans la part de mailles
    qui rapportent (F4).
  \item \textbf{Doublons :} l'identité d'une observation est (station, élément, date) ; c'est la clé primaire,
    un doublon ne peut donc pas être stocké.
  \item \textbf{Unités :} valeurs stockées en entiers bruts NOAA (dixièmes de °C ou de mm), convertis à l'affichage seulement.
\end{itemize}

\paragraph{Heure d'observation.} Le champ \cmd{OBS-TIME} est vide pour 100\,\% de nos lignes et 98\,\% des
valeurs proviennent de l'ECA\&D (\cmd{sflag} E), qui agrège des réseaux nationaux dont la « journée » de mesure
peut différer (en particulier pour la pluie, cumulée d'un matin au suivant). Une même date peut donc couvrir des
fenêtres de 24~h décalées entre stations. Conséquence : une comparaison entre stations au jour près peut être
décalée d'un jour. F4 en tient compte en regroupant les jours consécutifs en un seul événement ; F2 et F3
ne comparent pas des stations au jour près.

% ======================================================================================================================
\section{Services}

Interface en ligne de commande (\cmd{uv run python -m weatherscope.cli}), identique en local et en distribué.
Chaque résultat contient un champ \cmd{sources} que la commande \cmd{trace} transforme en liste des observations
d'origine (station, date, élément, valeur, flags, lot). L'option \cmd{-{}-json} donne le résultat complet
en ordre canonique.

\begin{description}
  \item[F1 -- Explorer une station.] Entrées : zone ou rectangle, nom ; puis station et intervalle.
    Sortie : par élément et par mois, jours attendus, présents, valides, invalides, traces, disponibilité,
    moyenne (températures) ou cumul (pluie), min, max.
    \emph{Exemple :} Madrid Retiro (SP000003195), juillet 2022 : TMAX moyenne 36,78~°C (max 40,7~°C),
    29 jours valides sur 31 (disponibilité 0,94).

  \item[F2 -- Épisodes prolongés.] Entrées : zone, période, élément, seuil, durée minimale.
    Un jour compte s'il a une valeur valide $\ge$ seuil ; un épisode est une suite maximale de dates consécutives
    qui comptent. Un jour absent ou invalide coupe l'épisode ; les dates étant continues, un épisode du 30 décembre
    au 2 janvier reste un seul épisode. Sortie : station, début, fin, durée, dépassement moyen, maximum.
    \emph{Exemple :} Espagne 2015--2024, TMAX $\ge$ 40~°C pendant 3 jours ou plus : 220 épisodes ; le plus long à
    SPE00120503, du 7 au 17 juillet 2022 (11 jours, max 45,2~°C). Pour PRCP $\ge$ 5~mm pendant 3 jours en Allemagne, 114 épisodes
    chevauchent le Nouvel An et sont bien comptés une seule fois.

  \item[F3 -- Mois anormaux.] Entrées : zones, années, élément. Référence : 1991--2020 (30 ans).
    Un mois-station est éligible s'il a au moins 80\,\% de jours valides ; une normale station $\times$ mois exige
    au moins 24 années éligibles sur 30. L'anomalie d'une zone est la moyenne des anomalies par maille de 0,5°,
    puis la moyenne des mailles, pour ne pas surpondérer les régions denses ; le nombre de mailles et de stations
    utilisées est renvoyé et un mois doit couvrir au moins 5 mailles. Classement par |anomalie|.
    \emph{Exemple :} TMAX, 2015--2024 : décembre 2015 en Allemagne (+5,16~°C, 176 mailles, 313 stations),
    avril 2023 en Espagne (+3,61~°C).

  \item[F4 -- Événements régionaux.] Entrées : zone, période, élément, quantile (95\,\%), seuils de surface.
    Un jour est inhabituel pour une station si sa valeur dépasse son quantile 95\,\% du même mois sur 1991--2020.
    Pour tenir compte de la densité, on compte des mailles et non des stations : un jour est régional si au moins
    3 mailles et 25\,\% des mailles qui rapportent sont touchées ; les jours régionaux consécutifs forment un
    événement, de sévérité égale à la somme des parts de mailles touchées. Un jour où une seule maille est touchée
    est un extrême isolé. Preuves : chaque observation inhabituelle avec son seuil.
    \emph{Exemple :} Allemagne, TMAX, 2015--2024 : 166 événements ; le plus sévère du 1er au 12 novembre 2015
    (98\,\% des mailles le 8). 185 jours d'extrême isolé, souvent à la Zugspitze (GM000004155, 2\,964~m).

  \item[F5 -- Mise à jour.] Entrée : fichiers NOAA réels (\cmd{by\_station} ou \cmd{by\_year}).
    Identité : (station, élément, date). Politique : une valeur ou un flag différent \emph{écrase} l'ancienne
    version, qui est \emph{archivée} dans \cmd{observations\_history} avec son lot. Une station allemande ou
    espagnole inconnue est ajoutée d'après \cmd{ghcnd-stations.txt}, les autres pays sont ignorés ; un lot contenant
    deux valeurs différentes pour une même observation est refusé. Les agrégats mensuels des mois touchés
    sont recalculés dans la même transaction.
    \emph{Exemple :} fichier \cmd{by\_year/2025} complet : 697\,993 lignes importées en 35~s, 18,9~M lignes
    d'autres pays ignorées ; second import : 0 nouvelle, 0 modifiée. Réimport des 21,46~M lignes : aucun changement.
\end{description}

% ======================================================================================================================
\section{Partie I -- Prototype local}

\paragraph{Technologie.} PostgreSQL 17.2, processus serveur séparé (conteneur Docker, port réseau).
Nos données sont tabulaires et de forme fixe, et F2--F4 sont des analyses (suites de jours, écarts à une moyenne,
croisement de stations aux mêmes dates) que les fonctions de fenêtre SQL expriment directement. Surtout, le même SQL
tourne sur Citus (Partie II), qui distribue PostgreSQL : un seul code pour les deux prototypes, ce qui facilite la
garantie de résultats identiques. Nous avons écarté MongoDB (F2--F4 en pipelines d'agrégation longs à écrire en deux jours),
Cassandra et ClickHouse (jamais pratiqués).

\paragraph{Organisation des données.}
\begin{itemize}
  \item \cmd{observations} : une ligne par mesure (station, élément, date, valeur entière, mflag, qflag, sflag, lot),
    clé primaire (station, élément, date). Ce format reprend exactement celui de la NOAA (chaque mesure a ses flags)
    et sert d'identité pour F5.
  \item \cmd{stations} : métadonnées et maille de 0,5° (\cmd{cell\_lat}, \cmd{cell\_lon}) utilisée par F3 et F4.
  \item \cmd{monthly\_stats} : par station, élément et mois, comptes (lignes, valides, invalides, traces, P),
    somme, min, max des valeurs valides. Maintenue par F5.
  \item \cmd{batches}, \cmd{observations\_history}, \cmd{observations\_staging} : provenance, historique et chargement de F5.
\end{itemize}

\paragraph{Accès et mécanismes.}
F1 et F2 lisent une station ou une zone sur une période : la clé primaire (station, élément, date) est aussi
l'index qui sert ces lectures. F2 utilise la méthode « date moins rang » : dans une suite de jours consécutifs,
\cmd{date - row\_number()} est constant, ce qui donne les épisodes en une seule passe triée par station.
F3 ne lit que \cmd{monthly\_stats} (86~Mo au lieu de 2,8~Go). Tous les calculs sont faits sur des entiers et
des \cmd{numeric}, arrondis seulement à la sortie : l'ordre des additions change en distribué et des flottants
donneraient des résultats légèrement différents. L'import (F5) passe par \cmd{COPY} vers une table de chargement puis
\cmd{INSERT \ldots{} ON CONFLICT DO UPDATE \ldots{} WHERE} la valeur change, ce qui rend l'import idempotent.
Taille en base : 2,9~Go (table 1,6~Go, index 1,3~Go). Chargement complet : 19~min.

\paragraph{Mesures.} \cmd{bench/partie1.py}, base complète, médiane de 5 exécutions après une exécution
d'échauffement (cache chaud) ; pour chaque comparaison, les deux versions renvoient exactement les mêmes lignes.

\paragraph{Décision justifiée par une mesure.} F3 (Allemagne et Espagne, TMAX, 2015--2024) prend 0,59~s sur
\cmd{monthly\_stats} contre 5,81~s en recalculant les mois depuis \cmd{observations} : les agrégats mensuels,
tenus à jour par F5, rendent F3 dix fois plus rapide pour 86~Mo de plus.

\paragraph{Conception alternative.} Une ligne par station et par jour, une colonne par élément avec ses flags
(\cmd{sql/alt\_wide.sql}), construite en 119~s à partir de \cmd{observations}.

\begin{table}[H]
\centering\small
\begin{tabular}{lrr}
\toprule
 & Une ligne par mesure (retenu) & Une ligne par jour (alternative) \\
\midrule
Lignes / taille avec index & 21,46~M / 2\,827~Mo & 7,91~M / 868~Mo \\
F1 : une station, 10 ans, 3 éléments & 0,010~s & 0,016~s \\
F2 : épisodes TMAX $\ge$ 40~°C, Espagne, 10 ans & 0,104~s & 0,214~s \\
TMAX $-$ TMIN $>$ 20~°C le même jour, Espagne, 10 ans & 2,04~s & 0,27~s \\
\bottomrule
\end{tabular}
\end{table}

La table large est 3,3 fois plus petite et 7 fois plus rapide quand une requête croise plusieurs éléments du même
jour. Mais nos services lisent un élément à la fois : la clé (station, élément, date) mène directement aux bonnes
lignes, alors que la table large lit toutes les colonnes de la station, d'où F1 et F2 1,5 à 2 fois plus lents.
Elle compliquerait aussi F5 (une ligne porte trois observations, avec trois jeux de flags et trois historiques).
Nous gardons donc une ligne par mesure ; la taille de la table large serait un argument si la base ne tenait plus
en mémoire (E1).

% ======================================================================================================================
\section{Partie II -- Prototype distribué}

\paragraph{Déploiement.} Citus 13.0 (PostgreSQL 17.2) : un coordinateur et 2 workers, conteneurs distincts avec
stockage et ports séparés. L'application ne connaît que le coordinateur : seule la variable \cmd{DATABASE\_URL} change.
\afaire{fichier de déploiement, ressources allouées.}

\paragraph{Placement.} \cmd{observations} est répartie par hachage de la station en 32 shards (16 par worker).
\cmd{observations\_history}, \cmd{observations\_staging} et \cmd{monthly\_stats} sont réparties de la même façon et
\emph{colocalisées} : les mêmes stations sont sur le même worker pour toutes ces tables. \cmd{stations} et
\cmd{batches} (774 lignes de métadonnées, pas des observations) sont des tables de référence copiées sur chaque worker.
Pourquoi la station : toutes les mesures d'une station sont au même endroit, donc F1 ne contacte qu'un worker,
les épisodes de F2 (fenêtres découpées par station) et F5 se calculent entièrement sur chaque worker.
Écartés : la date (F1 et F2 toucheraient tous les workers et F2 devrait recoller des épisodes coupés) et le pays
(2 valeurs seulement, très déséquilibré).

\paragraph{Instances contactées et combinaison.} Une requête sur une station (F1) est routée vers un seul shard
(\cmd{Task Count: 1}). Une requête sur une zone (F2, F3, F4) s'exécute sur les 32 shards ; chaque worker renvoie
des résultats partiels (épisodes complets pour F2, sommes et comptes pour F3 et F4) que le coordinateur fusionne.
Sur un cluster de test (65 stations), \cmd{-{}-explain} montre pour F1 \cmd{Task Count: 1} vers un seul worker
avec un parcours de l'index de la clé primaire, et pour F3 32 tâches réparties sur les deux workers, dont le résultat
est fusionné par le coordinateur. \afaire{extraits sur la base complète.}

\paragraph{Équilibre et écritures.} \afaire{répartition des lignes et de la taille par worker (\cmd{citus\_shards}),
temps du chargement distribué.}

\paragraph{Résultats identiques.} Toutes les sorties \cmd{-{}-json} sont en ordre canonique (égalités départagées
par station, date ou zone). Sur le même sous-ensemble chargé en local et sur Citus (65 stations, 1,83~M observations),
12 requêtes couvrant tous les services, dont F2--F4 sur des zones réparties sur les deux workers, donnent des sorties
identiques octet pour octet. \afaire{comparaison automatique sur les données complètes.}

\paragraph{Configuration alternative.} 3 workers (\cmd{citus\_add\_node} puis rééquilibrage), même code.
\afaire{mesures.}

\paragraph{Perte d'une instance.} Les shards ne sont pas répliqués. Si un worker est arrêté, toute requête qui en a
besoin échoue côté base, et l'interface affiche « ÉCHEC » avec la cause, sans aucun résultat partiel (code 2).
Essai sur le cluster de test, worker w2 arrêté : F1 sur GM000001474 (stockée sur w2) renvoie
« ÉCHEC : la base ou une de ses instances est indisponible \ldots{} connection to the remote node w2 failed \ldots{}
Aucun résultat partiel n'est affiché » ; F1 sur une station de w1 répond normalement ; F3 sur l'Allemagne échoue
de la même façon, puisqu'il lui faut les deux workers. Rien n'est récupéré : les données de w2 sont indisponibles
jusqu'à son redémarrage.
Un système en production répliquerait chaque shard (ou chaque worker) pour continuer à répondre ;
notre déploiement sur un seul poste ne le fait pas et n'est pas équivalent à un cluster de machines indépendantes.

% ======================================================================================================================
\section{Expériences}

Poste : processeur 8 cœurs / 16 threads, 13~Go de RAM, SSD NVMe ; PostgreSQL 17.2, Citus 13.0, \cmd{shared\_buffers}
2~Go en local. Toutes les instances partagent ce processeur, cette mémoire et ce disque.

\paragraph{E1 -- Volume.} \afaire{3 sous-ensembles (par ex. 2015--2024 = 7,14~M, 2005--2024, 1991--2024 = 21,46~M) :
temps d'import, temps des services, taille. La base complète (2,9~Go) tient-elle en mémoire ?}

\paragraph{E2 -- Local contre distribué.} \afaire{local, Citus-2, Citus-3 sur une requête étroite (F1),
une agrégation longue (F3) et une analyse multi-stations (F4) ; médiane et dispersion, conditions de cache.}

\paragraph{E3 -- Charge variable.} Prédiction écrite avant les mesures (\cmd{docs/e3\_prediction.md}) :
le local l'emporte sur la routine (F1 : surcoût fixe du coordinateur) ; Citus-2 l'emporte sur les études régionales
grâce au parallélisme de F2 ; dans un mélange, le basculement se fait vers 1 à 5\,\% de requêtes régionales.
\afaire{mesures pour 0, 1, 5, 25, 50, 100\,\% de requêtes régionales, et confrontation avec la prédiction.}

\paragraph{E4 -- Limite (bonus).} Piste : F4 recalcule les quantiles de référence (30 ans) à chaque requête
(29~s sur l'Allemagne) ; une table de seuils tenue à jour par F5 les rendrait immédiats, au prix d'écritures en plus.
\afaire{si mesuré.}

% ======================================================================================================================
\section{Vérification indépendante}

\cmd{verify/} relit directement les fichiers NOAA en Python, sans la base ni le code des services,
avec des calculs exacts (fractions). \cmd{verify/compare.py} compare automatiquement ses résultats à ceux des
services, après tri canonique : F2 TMAX Espagne (220 épisodes), F2 PRCP Allemagne (16\,126 épisodes),
F3 TMAX Allemagne et F3 PRCP Espagne (120 mois chacun, classement compris) sont identiques.
Seul écart rencontré pendant le développement : l'arrondi, que Python fait au pair (13,125 $\to$ 13,12) et PostgreSQL
loin de zéro (13,13) ; les deux côtés arrondissent désormais comme PostgreSQL.

% ======================================================================================================================
\section{Limites}

\begin{itemize}
  \item Heure d'observation inconnue : comparaisons entre stations au jour près approximatives (section 2).
  \item F4 utilise un seuil fixe (quantile 95\,\%) et des mailles de 0,5° : un relief marqué (Zugspitze) produit
    des extrêmes isolés qui sont réels mais peu représentatifs de la région.
  \item Les normales exigent 24 années sur 30 : les stations récentes ne participent pas à F3 et F4.
  \item 9 stations espagnoles ont une longitude de signe inversé dans les métadonnées NOAA (annexe~\ref{ann:anomalies}) :
    elles sont rangées dans une mauvaise maille pour F3 et F4. Nous ne corrigeons pas les métadonnées sources.
  \item Le distribué tourne sur un seul poste, sans réplication.
\end{itemize}

% ======================================================================================================================
\appendix
\renewcommand{\thesection}{\Alph{section}}
\section{Anomalies présentes dans les données}\label{ann:anomalies}

Calculées sur les fichiers NOAA (\cmd{verify/data\_report.py}, détail dans \cmd{docs/data\_report.md}),
1991--2024 ; valeurs en dixièmes comme dans les fichiers.

\begin{table}[H]
\centering\small
\begin{tabularx}{\textwidth}{>{\raggedright\arraybackslash}p{4.2cm}r>{\raggedright\arraybackslash}X}
\toprule
Anomalie & Fréquence & Exemple (station, date, valeur) \\
\midrule
\cmd{qflag} I (incohérence interne) & 682 & GM000003218, 1997-01-01, TMAX = $-151$ \\
\cmd{qflag} O (valeur aberrante) & 320 & GME00121150, 2017-06-29, PRCP = 1969 \\
\cmd{qflag} N (zéro suspect) & 192 & GME00123382, 2006-04-17, TMAX = 0 \\
\cmd{qflag} S (incohérence spatiale) & 172 & GME00102348, 1993-02-04, TMIN = 51 \\
\cmd{qflag} K (série répétée) & 26 & GME00132550, 1998-08-27, PRCP = 1 \\
Pluie trace (\cmd{mflag} T) & 267 & SPW00013025, 1991-01-08, PRCP = 0 \\
Pluie > 200~mm sans \cmd{qflag} & 15 & SP000007038, 2019-09-13, PRCP = 2037 \\
Station qui s'arrête pendant 2015--2024 & 131 & GME00102236, dernière observation 2019-04-07 \\
Station qui commence pendant 2015--2024 & 32 & GME00121486, première observation 2021-06-02 \\
Trou de plus de 30 jours (2015--2024) & 94 & GME00102169, du 2015-01-13 au 2015-03-19 \\
Longitude de signe inversé (\cmd{ghcnd-stations.txt}) & 9 stations & SPE00101035 (Alicante El Altet) : +0,5706, alors que la même station sous SPM00008360 est à $-0{,}558$ \\
\bottomrule
\end{tabularx}
\end{table}

Recherchées mais absentes : paire TMAX < TMIN le même jour parmi les observations sans \cmd{qflag},
et pluie \cmd{mflag} P ; les règles de traitement correspondantes restent implémentées.

\end{document}
